% from ega4-9-fr.tex
\oldpage[IV-3]{76}

Alors, dans les propositions (9.6.1) et (9.6.2, (I)), on peut partout remplacer
$f_s$ par sa restriction $f_s|Z_s:Z_s\to Y_s$ sans changer les conclusions. En
effet, la vérification de (9.2.1, (i)) se fait comme dans (9.6.1) et (9.6.2).
D'autre part, dans la réduction au cas où $S$ est noethérien, faite dans
(9.2.3), si $Z=q^{-1}(Z_0)$, où $q:X\to X_0$ est la projection canonique et
$Z_0$ une partie constructible de $X_0$ (8.3.11), le fait que
$(Z_0)_{s_0}$ soit ouvert dans $(X_0)_{s_0}$, pour $s_0=p(s)$, découle de
(2.4.10) et de ce que la projection $X_s\to(X_0)_{s_0}$ est surjective. On est
donc ramené à vérifier (9.2.1, (ii)) sous les nouvelles hypothèses. Or, comme
$Z_\eta$ est ouvert dans $X_\eta$, il existe un ouvert $Z'\subset X$ tel que
$Z_\eta=Z'\cap X_\eta$~; comme $X$ est alors noethérien, $Z'$ est
constructible, et il en est de même de $Z$ par hypothèse~; on conclut donc de
(9.5.2) et de ($\mathbf 0_{\mathrm{III}}$, 9.2.2) qu'il y a un voisinage $U$
de $\eta$ dans $S$ tel que $Z_s=Z'_s$ pour $s\in U$. Remplaçant $S$ par $U$,
on peut donc se borner au cas où $Z$ est ouvert dans $X$, et alors on est
ramené à ce qui a été prouvé dans (9.6.1) et (9.6.2).
\end{remarks}

\subsection{Constructibilité des propriétés de séparabilité, d'irréductibilité
géométrique et de connexité géométrique}
\label{subsection:IV.9.7-fr}

\begin{lemma}[9.7.1]
\label{IV.9.7.1-fr}
Soient $S$ un préschéma irréductible de point générique $\eta$, $f:X\to S$
un morphisme de présentation finie. Si $X_\eta$ a $n$ composantes irréductibles
(resp. $n'$ composantes connexes), il existe un voisinage ouvert $U$ de $\eta$
dans $S$ tel que pour tout $s\in U$, $X_s$ ait au moins $n$ composantes
irréductibles (resp. $n'$ composantes connexes).
\end{lemma}

On peut se borner au cas où $S$ est affine, donc $X$ quasi-compact et
quasi-séparé. Soient $V_i$ ($1\leqslant i\leqslant n$) les intérieurs des
composantes irréductibles de $X_\eta$, qui sont deux à deux disjoints et
quasi-compacts, et dont la réunion est dense dans $X_\eta$~; en vertu de
(8.2.11), appliqué suivant la méthode de (8.1.2, \emph{a)}), il existe pour
chaque $i$ un ouvert quasi-compact $W_i$ dans $X$ tel que
$W_i\cap X_\eta=V_i$~; comme $X$ est quasi-séparé, les intersections
$W_i\cap W_j$ sont quasi-compactes (1.2.7), donc, en remplaçant $S$
par un voisinage de $\eta$, on peut supposer que $W_i\cap W_j=\varnothing$ pour
$i\neq j$ en vertu de (8.3.3). En outre, comme les $W_i$ sont constructibles,
il y a un voisinage $U$ de $\eta$ dans $S$ tel que pour tout $s\in U$ les
$(W_i)_s$ soient non vides (9.5.1 et 9.2.3) et que la réunion des $(W_i)_s$ soit
\emph{dense} dans $X_s$ (9.5.3 et 9.2.3). Cela étant, les composantes
irréductibles (en nombre fini) des $(W_i)_s$ sont aussi les composantes
irréductibles de la réunion des $(W_i)_s$ $(\mathbf 0_{\mathrm I},2.1.7)$, donc
les adhérences dans $X_s$ de ces composantes sont les composantes
irréductibles de $X_s$ $(\mathbf 0_{\mathrm I},2.1.6)$ et leur nombre est
évidemment $\geqslant n$.

Soient maintenant $C_j$ ($1\leqslant j\leqslant n'$) les composantes connexes
de $X_\eta$, qui sont des ouverts quasi-compacts de $X_\eta$, deux à deux
disjoints. Si on remplace les $V_i$ par les $C_j$ dans le raisonnement
précédent, on voit (utilisant deux fois (8.3.3)) qu'on peut supposer $X$ réunion
de $n'$ ouverts quasi-compacts $M_j$ ($1\leqslant j\leqslant n'$) deux à deux
disjoints et que, dans un voisinage de $\eta$, les $(M_j)_s$ soient non vides.
Comme la réunion des $(M_j)_s$ est $X_s$, les composantes connexes des $(M_j)_s$
sont les composantes connexes de $X_s$, donc leur nombre est $\geqslant n'$.

